\documentclass{resume}

\usepackage[left=0.4in,top=0.4in,right=0.4in,bottom=0.4in]{geometry}
\usepackage{fontawesome5}
\usepackage{enumitem}
\newcommand{\tab}[1]{\hspace{.2667\textwidth}\rlap{#1}} 
\newcommand{\itab}[1]{\hspace{0em}\rlap{#1}}
\newcommand{\CPP}{C\nolinebreak[4]\hspace{-.05em}\raisebox{.25ex}{\footnotesize\bf ++}}

\renewcommand{\sectionskip}{\smallskip}
\setlist[itemize]{itemsep=2.5pt, parsep=0pt, topsep=1.5pt, partopsep=0pt}

% Contact Info
\name{Sanchit Dass}
\address{0123456789 \\ Bengaluru \\ \href{mailto:mail@example.com}{mail@example.com}
 \\ \faGithub\ \href{https://github.com/example}{github.com/example} \\
   \faLinkedin\ \href{https://linkedin.com/in/example}{linkedin.com/in/example}
}

\begin{document}

    % 1. Professional Summary
    % \begin{rSection}{Professional Summary}
    %     Full Stack Software Engineer with experience building LLM-driven platforms, AI tooling, and scalable web
    %     applications across startups and enterprise environments. Strong background in Python, JavaScript
    %     and React with proven delivery of production AI pipelines, data extraction platforms, and cross-platform
    %     applications.
    % \end{rSection}

    % 2. Work Experience
    \begin{rSection}{Work Experience}
        \textbf{Applied AI Engineer} \hfill Feb 2025 --- Present \\
        Tario AI \hfill \textit{Bengaluru, India}
        \begin{itemize}
            \item Building an agentic SDR system to automate lead research, enrichment, and outreach using LangChain and LangGraph.
            \item Designed FastAPI-based microservices including an SDR copilot chatbot and a batched email generation system using message queues.
            \item Built pipelines combining LinkedIn + web data with LLM agents to generate context aware personalized emails at scale.
            \item Improved cost and latency of LLM workflows through prompt optimization, batching, and system-level changes.
            \item Contributed across backend and frontend to ship production features end-to-end.
            \item Maintaining a production CRM backend built entirely in Spring Boot, supporting core sales workflows across 3+ internal teams.
        \end{itemize}

        \textbf{Software Engineer Specialist} \hfill Jul 2025 --- Nov 2025 \\
        xAI \hfill \textit{Remote, India}
        \begin{itemize}
            \item Built dataset curation and validation tooling across Python, JS, and \CPP\ to improve coding model performance.
            \item Investigated model failures using execution traces and prompt analysis, identifying root causes for fixes.
            \item Reduced evaluation runtime by 50\% via parallelized execution pipelines.
            \item Validated AI-generated code across real-world projects to ensure correctness and performance.
        \end{itemize}

        \textbf{Founding Engineer} \hfill Jan 2025 --- Apr 2025 \\
        FetchFox \hfill \textit{Remote, USA}
        \begin{itemize}
            \item Built and scaled a data extraction platform (Next.js + Python) used by 1,000+ users.
            \item Implemented integrations with Snowflake and Google Sheets using APIs and OAuth.
            \item Reduced LLM costs by 25\% through preprocessing and input optimization.
            \item Built internal benchmarking tools to evaluate scraping quality across models.
        \end{itemize}

        \textbf{Intern} \hfill Apr 2024 --- Oct 2024 \\
        ICPC Foundation \hfill \textit{Remote, USA}
        \begin{itemize}
            \item Built a Flutter-based social platform (auth, feeds, messaging) for competitive programmers.
            \item Designed real-time notification systems used during ICPC World Finals.
        \end{itemize}

        \textbf{Software Engineer} \hfill Jul 2020 --- Jun 2022 \\
        National Instruments \hfill \textit{Bengaluru, India}
        \begin{itemize}
            \item Developed and maintained production-grade \CPP\ APIs for RFmx signal analysis software, contributing to 4 major product releases.
            \item Led architectural changes enabling side by side RFmx installations, ensuring backward compatibility across releases by redesigning installation and versioning workflows.
            \item Reduced CI build time by 27\% by optimizing Jenkins pipelines, directly improving developer iteration speed.
            \item Built an automated Python-based testing framework integrated into nightly builds, significantly improving bug detection, triage speed, and release confidence.
        \end{itemize}
    \end{rSection}

    % 3. Skills
    \begin{rSection}{Skills}
        \begin{tabular}{ @{} >{\bfseries}l @{\hspace{6ex}} l }
            Programming & Python, JavaScript, Java \\
            AI / LLM & LangChain, LangGraph, RAG, Prompt Engineering \\
            Backend & FastAPI, Spring Boot, Node.js, Django \\
            Frontend & React, Next.js \\
            Systems & Docker, Message Queues \\
            Cloud & Amazon Web Services \\
        \end{tabular}\\
    \end{rSection}

    % 4. Education
    \begin{rSection}{Education}
        \textbf{MS, Computer Science}, Purdue University \hfill {Aug 2022 --- Dec 2023}\\
        % \begin{itemize}
        %     \item GPA (Only include if higher than 3.5 / 4.0), Class of Honors (Only include if good)
        %     \item Other school awards like Dean’s List, scholarships, community awards, etc.
        %     \item Relevant Coursework: Advanced Class A, Advanced Class B, Advanced Class C
        % \end{itemize}
        \textbf{BE, Computer Science}, BMS College of Engineering \hfill {Aug 2016 --- Sep 2020}
    \end{rSection}

    % 5. Projects
    % \begin{rSection}{Projects}

    % \textbf{Role}, {Project Name} (\href{https://myproject.com}{myproject.com}) \hfill {MMM YYYY --- MMM YYYY}
    % \begin{itemize}
    %     \item Action verb + what you did + reason, outcome or quantified results. Talk about technologies used.
    %     \item Action verb + what you did + reason, outcome or quantified results. Talk about technologies used.
    % \end{itemize}

    % \textbf{Role}, {Project Name} (\href{https://github.com/my-project}{github.com/my-project}) \hfill {MMM YYYY --- MMM YYYY}
    % \begin{itemize}
    %     \item Action verb + what you did + reason, outcome or quantified results. Talk about technologies used.
    %     \item Action verb + what you did + reason, outcome or quantified results. Talk about technologies used.
    % \end{itemize}

    % \end{rSection}

    % 6. Awards
    % \begin{rSection}{Awards}
    %     \begin{itemize}
    %         \item \textbf{MIT Women's Health AI Hackathon} --- Winners of Most Innovative Idea \hfill {May 2025}
    %     \end{itemize}
    % \end{rSection}
\end{document}
